% ============================================================================
% Database Systems Lab — Week 04: JOINs & Relational Data Modelling
% Lab Report
% ============================================================================
\documentclass[12pt,a4paper]{article}

% ============================================================================
% PACKAGES
% ============================================================================
\usepackage[utf8]{inputenc}
\usepackage[T1]{fontenc}
\usepackage{geometry}
\usepackage{graphicx}
\usepackage{xcolor}
\usepackage{listings}
\usepackage{hyperref}
\usepackage{booktabs}
\usepackage{titlesec}
\usepackage{fancyhdr}
\usepackage{enumitem}
\usepackage{amsmath}
\usepackage{float}
\usepackage{caption}
\usepackage{subcaption}
\usepackage{tcolorbox}
\usepackage{pgfplots}
\pgfplotsset{compat=1.18}

% ============================================================================
% PAGE SETUP
% ============================================================================
\geometry{top=2.5cm,bottom=2.5cm,left=2.5cm,right=2.5cm}
\pagestyle{fancy}
\fancyhf{}
\rhead{\thepage}
\lhead{Database Systems Lab — Week 04}
\rfoot{JOINs \& Relational Data Modelling}
\renewcommand{\headrulewidth}{0.4pt}
\renewcommand{\footrulewidth}{0.4pt}

% ============================================================================
% COLOR DEFINITIONS
% ============================================================================
\definecolor{codegreen}{rgb}{0,0.6,0}
\definecolor{codegray}{rgb}{0.5,0.5,0.5}
\definecolor{codepurple}{rgb}{0.58,0,0.82}
\definecolor{backcolour}{rgb}{0.95,0.95,0.92}
\definecolor{primaryblue}{RGB}{0,102,204}
\definecolor{secondarygray}{RGB}{80,80,80}

% ============================================================================
% SQL LISTING STYLE
% ============================================================================
\lstdefinestyle{sqlstyle}{
    language=SQL,
    backgroundcolor=\color{backcolour},
    commentstyle=\color{codegreen}\itshape,
    keywordstyle=\color{primaryblue}\bfseries,
    numberstyle=\tiny\color{codegray},
    stringstyle=\color{codepurple},
    basicstyle=\ttfamily\footnotesize,
    breakatwhitespace=false,
    breaklines=true,
    captionpos=b,
    keepspaces=true,
    numbers=left,
    numbersep=5pt,
    showspaces=false,
    showstringspaces=false,
    showtabs=false,
    tabsize=2,
    frame=single,
    rulecolor=\color{primaryblue},
    framerule=0.5pt
}
\lstset{style=sqlstyle}

% ============================================================================
% SECTION FORMATTING
% ============================================================================
\titleformat{\section}
  {\Large\bfseries\color{primaryblue}}
  {\thesection.}{0.5em}
  {\titlerule[0.8pt]}

\titleformat{\subsection}
  {\large\bfseries\color{secondarygray}}
  {\thesubsection.}{0.5em}{}

\titleformat{\subsubsection}
  {\normalsize\bfseries}
  {\thesubsubsection.}{0.5em}{}

\titlespacing*{\section}{0pt}{3.5ex plus 1ex minus .2ex}{2.3ex plus .2ex}
\titlespacing*{\subsection}{0pt}{3.25ex plus 1ex minus .2ex}{1.5ex plus .2ex}

% ============================================================================
% HYPERREF SETUP
% ============================================================================
\hypersetup{
    colorlinks=true,
    linkcolor=primaryblue,
    filecolor=magenta,
    urlcolor=primaryblue,
    citecolor=primaryblue,
}

% ============================================================================
% CUSTOM BOXES
% ============================================================================
\newtcolorbox{querybox}[1][]{
    colback=blue!5!white,
    colframe=primaryblue,
    fonttitle=\bfseries,
    title=Query #1,
    boxrule=0.5pt,
    arc=3pt,
    left=6pt,
    right=6pt,
    top=6pt,
    bottom=6pt
}

\newtcolorbox{analysisbox}[1][]{
    colback=green!5!white,
    colframe=green!60!black,
    fonttitle=\bfseries,
    title=Analysis #1,
    boxrule=0.5pt,
    arc=3pt,
    left=6pt,
    right=6pt,
    top=6pt,
    bottom=6pt
}

\newtcolorbox{resultbox}{
    colback=gray!5!white,
    colframe=secondarygray,
    fonttitle=\bfseries,
    title=Results \& Interpretation,
    boxrule=0.5pt,
    arc=3pt,
    left=6pt,
    right=6pt,
    top=6pt,
    bottom=6pt
}

% ============================================================================
% TITLE PAGE
% ============================================================================
\title{
    \vspace{-1cm}
    \rule{\linewidth}{0.5mm}\\
    \vspace{0.5cm}
    {\Huge \bfseries Database Systems Lab — Week 04}\\
    \vspace{0.3cm}
    {\Large JOINs \& Relational Data Modelling}\\
    \vspace{0.3cm}
    \rule{\linewidth}{0.5mm}\\
    \vspace{1cm}
}

\author{
    \textbf{Student Name:} [Muhammad Ayyan Khan]\\
    \textbf{Course:} Database Systems Lab (Data Science Program)\\
    \textbf{Instructor:} Muhammad Usama Afridi\\
    \textbf{Institute:} Institute of Management Sciences\\
    \vspace{0.5cm}
    \textbf{Date:} \today
}

\date{}

% ============================================================================
% DOCUMENT START
% ============================================================================
\begin{document}

\maketitle
\thispagestyle{empty}

% ============================================================================
% TABLE OF CONTENTS
% ============================================================================
\tableofcontents
\newpage

% ============================================================================
% SECTION 1: 10 QUERIES
% ============================================================================
\section{Section 1 — 10 Queries (35 points)}

This section presents all 10 required queries with SQL code, result screenshots,
and business-level interpretation.

% ----------------------------------------------------------------------------
% QUERY 1
% ----------------------------------------------------------------------------
\subsection{Query 1: Orders with Customer Names (INNER JOIN)}

\begin{querybox}[1]
\textbf{Objective:} Show order details with customer name and city instead of just customer\_id.
\end{querybox}

\begin{lstlisting}[caption={Query 1: Orders with Customer Names}]
SELECT o.order_id,
       c.name        AS customer_name,
       c.city,
       o.order_date,
       o.status,
       o.total_amount
FROM orders o
INNER JOIN customers c ON o.customer_id = c.customer_id
ORDER BY o.order_date DESC;
\end{lstlisting}

\begin{figure}[H]
    \centering
    \includegraphics[width=0.9\textwidth]{pdf/Query1—Your-firs-join:-orders-with-customer-names.png}
    \caption{Query 1 Results: Orders with Customer Names}
    \label{fig:query1}
\end{figure}

\begin{resultbox}
\textbf{Interpretation:} This INNER JOIN connects the orders table to the customers
table, replacing opaque customer\_id values with readable customer names and cities.
The results show that orders are sorted by date (newest first), allowing business
users to quickly identify recent transactions and their associated customers. This
is the most common join pattern in analytics — connecting transaction tables to
lookup/dimension tables.
\end{resultbox}

% ----------------------------------------------------------------------------
% QUERY 2
% ----------------------------------------------------------------------------
\subsection{Query 2: Full Order Receipt (4-Table Join)}

\begin{querybox}[2]
\textbf{Objective:} Produce a complete order receipt showing customer name, order date,
product details, quantity, unit price, and line total.
\end{querybox}

\begin{lstlisting}[caption={Query 2: Full Order Receipt}]
SELECT c.name         AS customer,
       o.order_id,
       o.order_date,
       o.status,
       p.product_name,
       p.category,
       oi.quantity,
       oi.unit_price,
       oi.quantity * oi.unit_price AS line_total
FROM order_items oi
INNER JOIN orders o ON oi.order_id = o.order_id
INNER JOIN customers c ON o.customer_id = c.customer_id
INNER JOIN products p ON oi.product_id = p.product_id
ORDER BY o.order_id, p.category;
\end{lstlisting}

\begin{figure}[H]
    \centering
    \includegraphics[width=0.9\textwidth]{pdf/Query2---Full-order-receipt.png}
    \caption{Query 2 Results: Complete Order Receipt}
    \label{fig:query2}
\end{figure}

\begin{resultbox}
\textbf{Interpretation:} This four-table join represents the complete transaction
picture — the heart of any e-commerce analytics system. Every business report
("What did customer X buy?", "How much did category Y sell?") starts from this
exact join structure. The line\_total calculation (quantity × unit\_price) provides
the revenue contribution of each item, essential for profitability analysis.
\end{resultbox}

% ----------------------------------------------------------------------------
% QUERY 3
% ----------------------------------------------------------------------------
\subsection{Query 3: Customers Who Have Never Ordered (LEFT JOIN)}

\begin{querybox}[3]
\textbf{Objective:} Find customers who signed up but have never placed any order
(churn/re-engagement analysis).
\end{querybox}

\begin{lstlisting}[caption={Query 3: Customers Who Never Ordered}]
SELECT c.customer_id,
       c.name,
       c.email,
       c.city,
       c.signup_date
FROM customers c
LEFT JOIN orders o ON c.customer_id = o.customer_id
WHERE o.order_id IS NULL
ORDER BY c.signup_date;
\end{lstlisting}

\begin{figure}[H]
    \centering
    \includegraphics[width=0.9\textwidth]{pdf/Query3----Customers-who-have-never-ordered.png}
    \caption{Query 3 Results: Customers Who Never Ordered}
    \label{fig:query3}
\end{figure}

\begin{resultbox}
\textbf{Interpretation:} The LEFT JOIN + WHERE IS NULL pattern (called an "anti-join")
identifies customers with zero orders. These are critical for business: they represent
acquisition costs with no ROI. Marketing teams use this list for re-engagement campaigns.
The signup\_date ordering helps prioritize recent signups who may just need a nudge.
\end{resultbox}

% ----------------------------------------------------------------------------
% QUERY 4
% ----------------------------------------------------------------------------
\subsection{Query 4: Products Never Ordered (LEFT JOIN)}

\begin{querybox}[4]
\textbf{Objective:} Find products that have never appeared in order\_items (inventory analysis).
\end{querybox}

\begin{lstlisting}[caption={Query 4: Products Never Ordered}]
SELECT p.product_id,
       p.product_name,
       p.category,
       p.price,
       p.stock_qty
FROM products p
LEFT JOIN order_items oi ON p.product_id = oi.product_id
WHERE oi.item_id IS NULL
ORDER BY p.category;
\end{lstlisting}

\begin{figure}[H]
    \centering
    \includegraphics[width=0.9\textwidth]{pdf/Query4—Products-never-ordered(before).png}
    \caption{Query 4 Results (Before): No Unsold Products}
    \label{fig:query4-before}
\end{figure}

\subsubsection*{Verification: Adding a Test Product}

Since all products had been ordered, I added a test product to verify the query works:

\begin{lstlisting}[caption={Adding Test Product}]
INSERT INTO products (product_name, category, price, stock_qty)
VALUES ('Unsold Widget', 'Test Products', 99.99, 100);
\end{lstlisting}

\begin{figure}[H]
    \centering
    \includegraphics[width=0.9\textwidth]{pdf/Query4--Products-never-ordered(after)EXplain analyze.png}
    \caption{Query 4 Results (After): Test Product Appears as Unsold}
    \label{fig:query4-after}
\end{figure}

\begin{resultbox}
\textbf{Interpretation:} Products that never sell tie up inventory budget and warehouse
space. In a real business, this query runs on a schedule to identify dead stock for
clearance or discontinuation. The EXPLAIN ANALYZE output shows PostgreSQL using a
Hash Left Join strategy, which is efficient for this type of anti-join pattern.
\end{resultbox}

% ----------------------------------------------------------------------------
% QUERY 5
% ----------------------------------------------------------------------------
\subsection{Query 5: Revenue by Customer with Join}

\begin{querybox}[5]
\textbf{Objective:} Show every customer's total orders and total spent, including
customers who have never ordered (showing 0).
\end{querybox}

\begin{lstlisting}[caption={Query 5: Revenue by Customer}]
SELECT c.name,
       c.city,
       COUNT(o.order_id)              AS total_orders,
       COALESCE(SUM(o.total_amount), 0) AS total_spent
FROM customers c
LEFT JOIN orders o ON c.customer_id = c.customer_id
GROUP BY c.customer_id, c.name, c.city
ORDER BY total_spent DESC;
\end{lstlisting}

\begin{figure}[H]
    \centering
    \includegraphics[width=0.9\textwidth]{pdf/Query5—Revenue-by-customer-with-join.png}
    \caption{Query 5 Results: Revenue by Customer}
    \label{fig:query5}
\end{figure}

\begin{resultbox}
\textbf{Interpretation:} This query combines aggregation (Week 3) with LEFT JOINs
(Week 4). The COALESCE function ensures customers with no orders show 0 instead of
NULL. Using COUNT(o.order\_id) instead of COUNT(*) correctly excludes NULL values
from non-matching rows. This is a classic customer lifetime value (CLV) calculation.
\end{resultbox}

% ----------------------------------------------------------------------------
% QUERY 6
% ----------------------------------------------------------------------------
\subsection{Query 6: Best-Selling Products}

\begin{querybox}[6]
\textbf{Objective:} Find the top 5 products by total units sold with revenue.
\end{querybox}

\begin{lstlisting}[caption={Query 6: Best-Selling Products}]
SELECT p.product_name,
       p.category,
       p.price,
       SUM(oi.quantity)                AS units_sold,
       SUM(oi.quantity * oi.unit_price) AS total_revenue
FROM products p
INNER JOIN order_items oi ON p.product_id = oi.product_id
GROUP BY p.product_id, p.product_name, p.category, p.price
ORDER BY units_sold DESC
LIMIT 5;
\end{lstlisting}

\begin{figure}[H]
    \centering
    \includegraphics[width=0.9\textwidth]{pdf/Query6—Best-selling-products.png}
    \caption{Query 6 Results: Top 5 Best-Selling Products}
    \label{fig:query6}
\end{figure}

\begin{resultbox}
\textbf{Interpretation:} This identifies the most popular products by volume and
revenue. The GROUP BY includes all non-aggregated columns (required in PostgreSQL).
Business use: inventory planning (stock more of top sellers), marketing (feature
these products), and supplier negotiations (leverage high-volume purchases).
\end{resultbox}

% ----------------------------------------------------------------------------
% QUERY 7
% ----------------------------------------------------------------------------
\subsection{Query 7: City-Level Sales Report}

\begin{querybox}[7]
\textbf{Objective:} For each city, show customer count, order count, total revenue,
and average order value.
\end{querybox}

\begin{lstlisting}[caption={Query 7: City-Level Sales Report}]
SELECT c.city,
       COUNT(DISTINCT c.customer_id) AS num_customers,
       COUNT(DISTINCT o.order_id)    AS num_orders,
       SUM(o.total_amount)           AS total_revenue,
       ROUND(AVG(o.total_amount), 2) AS avg_order_value
FROM customers c
INNER JOIN orders o ON c.customer_id = c.customer_id
GROUP BY c.city
ORDER BY total_revenue DESC;
\end{lstlisting}

\begin{figure}[H]
    \centering
    \includegraphics[width=0.9\textwidth]{pdf/Query7—City-level-sales-report.png}
    \caption{Query 7 Results: City-Level Sales Report}
    \label{fig:query7}
\end{figure}

\begin{resultbox}
\textbf{Interpretation:} Geographic analysis reveals market penetration and regional
preferences. COUNT(DISTINCT) prevents double-counting customers with multiple orders.
High-revenue cities may deserve targeted marketing or physical store expansion.
Low-revenue cities with many customers indicate low average spend — a pricing or
product mix issue.
\end{resultbox}

% ----------------------------------------------------------------------------
% QUERY 8
% ----------------------------------------------------------------------------
\subsection{Query 8: Order Details with Status Filter}

\begin{querybox}[8]
\textbf{Objective:} Show delivered orders from 2025, grouped by customer with
product diversity and spend metrics.
\end{querybox}

\begin{lstlisting}[caption={Query 8: Order Details with Status Filter}]
SELECT c.name                              AS customer,
       c.city,
       COUNT(DISTINCT p.product_id)        AS distinct_products_bought,
       SUM(oi.quantity)                    AS total_items,
       SUM(oi.quantity * oi.unit_price)    AS total_spent
FROM order_items oi
INNER JOIN orders o ON oi.order_id = o.order_id
INNER JOIN customers c ON o.customer_id = c.customer_id
INNER JOIN products p ON oi.product_id = p.product_id
WHERE o.status = 'delivered'
  AND o.order_date >= '2025-01-01'
GROUP BY c.customer_id, c.name, c.city
ORDER BY total_spent DESC;
\end{lstlisting}

\begin{figure}[H]
    \centering
    \includegraphics[width=0.9\textwidth]{pdf/Query8—Order-details-with-status-filter.png}
    \caption{Query 8 Results: Delivered Orders from 2025}
    \label{fig:query8}
\end{figure}

\begin{resultbox}
\textbf{Interpretation:} This combines joins, WHERE filters, GROUP BY, and aggregates
in one query. Filtering to 'delivered' status ensures we analyze completed transactions
only. The distinct\_products\_bought metric shows product diversity — high values
indicate customers exploring the catalog vs. repeat buyers of the same items.
\end{resultbox}

% ----------------------------------------------------------------------------
% QUERY 9
% ----------------------------------------------------------------------------
\subsection{Query 9: Session and Purchase Behaviour Together}

\begin{querybox}[9]
\textbf{Objective:} For every customer with sessions, show session count, orders,
total spent, and conversion status.
\end{querybox}

\begin{lstlisting}[caption={Query 9: Session and Purchase Behaviour}]
SELECT c.name,
       c.city,
       COUNT(DISTINCT s.session_id)       AS total_sessions,
       COUNT(DISTINCT o.order_id)         AS total_orders,
       COALESCE(SUM(DISTINCT o.total_amount), 0) AS total_spent,
       CASE
           WHEN COUNT(DISTINCT o.order_id) > 0 THEN 'Converted'
           ELSE 'Not Converted'
       END                                AS customer_status
FROM customers c
INNER JOIN user_sessions s ON c.customer_id = s.customer_id
LEFT JOIN orders o ON c.customer_id = o.customer_id
       AND o.status = 'delivered'
GROUP BY c.customer_id, c.name, c.city
ORDER BY total_sessions DESC;
\end{lstlisting}

\begin{figure}[H]
    \centering
    \includegraphics[width=0.9\textwidth]{pdf/Query9—Session-and-purchase-behaviour-together.png}
    \caption{Query 9 Results: Session and Purchase Behaviour}
    \label{fig:query9}
\end{figure}

\begin{resultbox}
\textbf{Interpretation:} This query connects browsing behaviour (sessions) to
purchases (orders). The filter on o.status is in the ON clause (not WHERE) to
preserve the LEFT JOIN behaviour. High sessions with no conversion indicates
UX/pricing problems. This is the foundation for conversion rate optimization.
\end{resultbox}

% ----------------------------------------------------------------------------
% QUERY 10
% ----------------------------------------------------------------------------
\subsection{Query 10: FULL OUTER JOIN — Data Completeness Audit}

\begin{querybox}[10]
\textbf{Objective:} Identify mismatches between customers and user\_sessions tables
(data quality audit).
\end{querybox}

\begin{lstlisting}[caption={Query 10: Data Completeness Audit}]
SELECT c.customer_id                     AS customer_table_id,
       c.name,
       s.customer_id                     AS session_table_cust_id,
       COUNT(s.session_id)               AS session_count,
       CASE
           WHEN c.customer_id IS NULL THEN 'Orphaned session — no customer record'
           WHEN s.customer_id IS NULL THEN 'Customer with no sessions'
           ELSE 'Matched'
       END                               AS audit_status
FROM customers c
FULL OUTER JOIN user_sessions s ON c.customer_id = s.customer_id
GROUP BY c.customer_id, c.name, s.customer_id
ORDER BY audit_status, c.name;
\end{lstlisting}

\begin{figure}[H]
    \centering
    \includegraphics[width=0.9\textwidth]{pdf/Query10—FULL-OUTER-JOIN:-data-completeness-audit.png}
    \caption{Query 10 Results: Data Completeness Audit}
    \label{fig:query10}
\end{figure}

\begin{resultbox}
\textbf{Interpretation:} FULL OUTER JOIN returns all rows from both tables, revealing
data integrity issues. "Orphaned sessions" indicate tracking bugs or deleted accounts.
"Customers with no sessions" may be legitimate (phone orders) or indicate tracking gaps.
This audit should run before any analysis to ensure data quality.
\end{resultbox}

% ============================================================================
% SECTION 2: ER DIAGRAM
% ============================================================================


% NOTE: If you have an ER diagram image, uncomment the figure below and add the image to pdf/ folder
% \begin{figure}[H]
%     \centering
%     \includegraphics[width=0.85\textwidth]{pdf/ER-Diagram.png}
%     \caption{Entity-Relationship Diagram: E-Commerce Schema}
%     \label{fig:er-diagram}
% \end{figure}

\subsection*{Entity-Relationship Diagram (Text Representation)}

\begin{tcolorbox}[colback=white,colframe=primaryblue,title=ER Diagram,coltitle=white]
\footnotesize
\begin{verbatim}
+------------------+       +------------------+       +------------------+
|    CUSTOMERS     |       |     ORDERS       |       |   ORDER_ITEMS    |
+------------------+       +------------------+       +------------------+
| PK customer_id   |<------| FK customer_id   |       | FK order_id      |
|    name          |  1:N  | PK order_id      |  1:N  | FK product_id    |
|    email         |       |    order_date    |------>|    quantity      |
|    city          |       |    status        |       |    unit_price    |
|    signup_date   |       |    total_amount  |       +------------------+
+------------------+       |    shipped_date  |                ^
         ^                 +------------------+                | 1:N
         |                                                    |
         | 1:N                                                |
         |                                          +------------------+
         |                                          |    PRODUCTS      |
         |                                          +------------------+
         |                                          | PK product_id    |
         +------------------------------------------|    product_name  |
                    1:N                             |    category      |
                                                    |    price         |
+------------------+                                |    stock_qty     |
|  USER_SESSIONS   |                                +------------------+
+------------------+
| PK session_id    |
| FK customer_id   |
|    session_date  |
|    device        |
|    converted     |
+------------------+
\end{verbatim}
\end{tcolorbox}

\textbf{Relationships:}
\begin{itemize}[leftmargin=*,noitemsep,topsep=3pt]
    \item \textbf{customers → orders:} One-to-Many (one customer can place many orders)
    \item \textbf{orders → order\_items:} One-to-Many (one order contains many items)
    \item \textbf{products → order\_items:} One-to-Many (one product can appear in many order items)
    \item \textbf{customers → user\_sessions:} One-to-Many (one customer can have many sessions)
\end{itemize}

\vspace{0.5cm}
\textbf{Primary Keys (PK):} customer\_id, order\_id, item\_id, product\_id, session\_id

\vspace{0.3cm}
\textbf{Foreign Keys (FK):}
\begin{itemize}[leftmargin=*,noitemsep,topsep=3pt]
    \item orders.customer\_id → customers.customer\_id
    \item order\_items.order\_id → orders.order\_id
    \item order\_items.product\_id → products.product\_id
    \item user\_sessions.customer\_id → customers.customer\_id
\end{itemize}

% ============================================================================
% SECTION 3: NORMALISATION REFLECTION
% ============================================================================
\section{Section 3 — Normalisation Reflection (15 points)}

\begin{tcolorbox}[
    colback=white,
    colframe=secondarygray,
    fonttitle=\bfseries,
    title=Option (b): Is the user\_sessions table in 3NF?,
    boxrule=0.5pt,
    arc=3pt,
    left=10pt,
    right=10pt,
    top=10pt,
    bottom=10pt
]

\subsection*{Analysis: user\_sessions Table Normalisation}

The \textbf{user\_sessions} table has the following structure:
\begin{itemize}[leftmargin=*,noitemsep,topsep=3pt]
    \item session\_id (PK)
    \item customer\_id (FK → customers.customer\_id)
    \item session\_date
    \item device
    \item converted
\end{itemize}

To determine if this table is in Third Normal Form (3NF), we verify three conditions:

\subsubsection*{1. First Normal Form (1NF) ✓}
The table satisfies 1NF because:
\begin{itemize}[leftmargin=*,noitemsep,topsep=2pt]
    \item Each column contains atomic (single) values
    \item There are no repeating groups
    \item Each row is uniquely identifiable by session\_id
    \item No column contains multiple values like "device1, device2" in one cell
\end{itemize}

\subsubsection*{2. Second Normal Form (2NF) ✓}
The table satisfies 2NF because:
\begin{itemize}[leftmargin=*,noitemsep,topsep=2pt]
    \item The primary key is a single column (session\_id), not composite
    \item All non-key columns (customer\_id, session\_date, device, converted) depend on the entire primary key
    \item There is no partial dependency issue since the key is not composite
\end{itemize}

\subsubsection*{3. Third Normal Form (3NF) ✓}
The table satisfies 3NF because:
\begin{itemize}[leftmargin=*,noitemsep,topsep=2pt]
    \item All non-key columns depend only on the primary key (session\_id)
    \item There are no transitive dependencies
    \item customer\_id is a foreign key, which is allowed in 3NF
    \item session\_date, device, and converted all describe the session itself, not some other non-key attribute
    \item No column can be derived from another non-key column
\end{itemize}

\subsection*{Conclusion}

\textbf{Yes, the user\_sessions table IS in Third Normal Form (3NF).}

\textbf{Justification:}
\begin{itemize}[leftmargin=*,noitemsep,topsep=3pt]
    \item Every attribute depends on the primary key (session\_id)
    \item The foreign key customer\_id correctly references the customers table
    \item There are no derived columns (e.g., we don't store customer\_name here)
    \item There are no transitive dependencies (e.g., device does not determine any other column)
\end{itemize}


\end{tcolorbox}

% ============================================================================
% SECTION 4: AI LEARNING LOG
% ============================================================================
\section{Section 4 — AI Learning Log (20 points)}

\subsection*{Entry 1: Debugging GROUP BY Error}

\textbf{Problem:} Query returned error: "column must appear in the GROUP BY clause
or be used in an aggregate function"

\textbf{Wrong Query:}
\begin{lstlisting}
SELECT p.product_name, p.category, p.price,
       SUM(oi.quantity) AS units_sold,
       SUM(oi.quantity * oi.unit_price) AS total_revenue
FROM products p
INNER JOIN order_items oi ON p.product_id = oi.product_id
ORDER BY units_sold DESC
LIMIT 5;
\end{lstlisting}

\textbf{AI Interaction:} Asked why the error occurred.

\textbf{Learning:} When using aggregate functions (SUM, COUNT, AVG), ALL non-aggregated
columns in SELECT must appear in GROUP BY. PostgreSQL is strict about this.

\textbf{Corrected Query:}
\begin{lstlisting}
SELECT p.product_name, p.category, p.price,
       SUM(oi.quantity) AS units_sold,
       SUM(oi.quantity * oi.unit_price) AS total_revenue
FROM products p
INNER JOIN order_items oi ON p.product_id = oi.product_id
GROUP BY p.product_id, p.product_name, p.category, p.price
ORDER BY units_sold DESC
LIMIT 5;
\end{lstlisting}

\subsection*{Entry 2: INNER JOIN vs LEFT JOIN Confusion}

\textbf{Problem:} Didn't understand when to use INNER vs LEFT JOIN for customer analysis.

\textbf{AI Interaction:} Asked for explanation of the difference with specific examples.

\textbf{Learning:}
\begin{itemize}[leftmargin=*,noitemsep]
    \item \textbf{INNER JOIN:} Only returns matching rows. Customers with no orders disappear.
    \item \textbf{LEFT JOIN:} Returns ALL left table rows. Non-matching right columns become NULL.
    \item Use LEFT JOIN when you need to include "zero" cases (customers who never ordered).
\end{itemize}

\subsection*{Entry 3: WHERE in ON Clause vs WHERE Clause}

\textbf{Problem:} Query 9 was excluding non-converted customers unexpectedly.

\textbf{Wrong Approach:}
\begin{lstlisting}
LEFT JOIN orders o ON c.customer_id = o.customer_id
WHERE o.status = 'delivered'  -- This turns LEFT JOIN into INNER JOIN!
\end{lstlisting}

\textbf{AI Interaction:} Asked why LEFT JOIN wasn't preserving all customers.

\textbf{Learning:} Filters on the RIGHT table of a LEFT JOIN must go in the ON clause,
not WHERE. WHERE filters AFTER the join, eliminating NULL rows you wanted to keep.

\textbf{Corrected:}
\begin{lstlisting}
LEFT JOIN orders o ON c.customer_id = o.customer_id
                    AND o.status = 'delivered'  -- Correct!
\end{lstlisting}

% ============================================================================
% SECTION 5: REFLECTION
% ============================================================================
\newpage
\section{Section 5 — Reflection (10 points)}

\subsection*{INNER JOIN vs LEFT JOIN: Understanding the Difference}

The fundamental difference between INNER JOIN and LEFT JOIN is about \textbf{inclusion}.
INNER JOIN is exclusive — it only returns rows where there's a match in both tables.
If a customer has never placed an order, an INNER JOIN between customers and orders
simply excludes that customer from the results. They become invisible. LEFT JOIN is
inclusive — it returns ALL rows from the left table regardless of whether there's a
match. Unmatched rows get NULL values for the right table's columns, but they appear.

I would use INNER JOIN when I only care about existing relationships: "Show me all
orders with their customers" (I don't need orphaned orders). I would use LEFT JOIN
when I need to find gaps or include zero cases: "Show me ALL customers with their
order counts" (including customers with zero orders).

\subsection*{Most Surprising Discovery}

The most surprising thing I discovered was how much insight is hidden in the
\textbf{absence} of relationships. Looking at individual tables, I could see
customers existed and orders existed. But only by joining them with LEFT JOIN
and filtering for NULLs did I discover that some customers had never ordered.
Similarly, products that had never been ordered only became visible through the
anti-join pattern. This changed my perspective: the most valuable business insights
often come from finding what's \textit{not} there — customers who churned, products
that don't sell, sessions that didn't convert. JOINs are not just about connecting
data; they're about revealing the gaps that drive business decisions.

The FULL OUTER JOIN audit was particularly eye-opening. I had assumed the database
would be perfectly consistent, but the possibility of orphaned records (sessions
without customers, customers without sessions) revealed that real-world data
pipelines have bugs. Data quality is not guaranteed — it must be audited.

% ============================================================================
% APPENDIX: PERFORMANCE ANALYSIS
% ============================================================================
\newpage
\section*{Appendix: Performance Analysis (EXPLAIN ANALYZE)}

\subsection*{Query 2 Execution Plan}

\begin{lstlisting}[caption={EXPLAIN ANALYZE for Query 2}]
EXPLAIN ANALYZE
SELECT c.name, o.order_id, o.order_date, o.status,
       p.product_name, p.category, oi.quantity, oi.unit_price,
       oi.quantity * oi.unit_price AS line_total
FROM order_items oi
INNER JOIN orders o ON oi.order_id = o.order_id
INNER JOIN customers c ON o.customer_id = c.customer_id
INNER JOIN products p ON oi.product_id = p.product_id
ORDER BY o.order_id, p.category;
\end{lstlisting}

\begin{figure}[H]
    \centering
    \includegraphics[width=0.9\textwidth]{pdf/EXPLAIN-ANALYZE.png}
    \caption{Query 2 Execution Plan}
    \label{fig:explain-analyze}
\end{figure}

\textbf{Operations Observed:}
\begin{itemize}[leftmargin=*,noitemsep]
    \item \textbf{Sort:} PostgreSQL sorts the final result by order\_id and category.
    \item \textbf{Hash Join:} Builds a hash table from one side and probes it for matches.
    Much faster than Nested Loop for large datasets.
    \item \textbf{Seq Scan:} Sequential scan reads entire tables (acceptable for small data).
\end{itemize}

On this small dataset, PostgreSQL chooses efficient hash joins. On millions of rows,
it would still use hash joins but might add parallel execution and index scans.

% ============================================================================
% BIBLIOGRAPHY
% ============================================================================
\newpage
\begin{thebibliography}{9}
\bibitem{week4manual}
Afridi, M. U. (2025). \textit{Database Systems Lab — Week 04 Manual: JOINs \& Relational Data Modelling}. Institute of Management Sciences.

\end{thebibliography}

\end{document}
